\chapter{Reading an unfamiliar proof one piece at a time}\label{ch:bridge}
\goal{Read a proof whose background you only partly have, without first studying a whole subject. The example is a proof that every sequence in a complete, totally bounded metric space has a convergent subsequence. Draw a map of what this one proof needs, learn only the pieces that are new to you, and rebuild the whole argument. Then use the same method to begin reading a formalized research proof, keeping in mind that its short final file rests on many other files, which hold most of the mathematics.}

\section{A new theorem on a familiar foundation}\label{ch17:sec:start}
Imagine meeting the following sentence in a research paper:
\begin{quote}
Every sequence in a complete, totally bounded metric space has a convergent subsequence.
\end{quote}
Most of its words are familiar: Chapter~\ref{ch:limits} defined metric spaces, convergence, Cauchy sequences and completeness, and Chapter~\ref{ch:contraction} used them all. Two expressions are new: ``totally bounded'' and ``subsequence.''

The conclusion is a form of what mathematicians call \emph{compactness}, which topology studies in great generality. Instead of first taking a course on that theory, ask what the two new ideas do in this particular argument, and learn exactly that much. A first exchange with an AI assistant might go like this.

\textbf{Reader:} I know metric spaces, Cauchy sequences and completeness. Describe how the proof of this statement works, without quoting any theorem I may not know.

\textbf{Assistant:} Keep infinitely many terms of the sequence inside a small region, then inside a smaller one, and so on. Pick later and later terms from these nested selections. The picked terms form a Cauchy sequence, and completeness supplies their limit.

\textbf{Reader:} Which facts do those steps rely on that I may not have met yet?

\textbf{Assistant:} Three. You need to know that the space can be covered by finitely many small balls; that if infinitely many indices are split into finitely many groups, some group still contains infinitely many of them; and how to choose indices that keep increasing.

The reader's side of this exchange is a \emph{minimal-background request}: it says what the reader already knows and asks only for what one particular proof needs. The words ``those steps'' pin down which proof, and that matters: a different proof of the same theorem could need different concepts. A good request also asks to see, right after each new concept, the step of the proof that uses it, which keeps the conversation tied to one argument instead of drifting into a survey of a subject.

You need not follow every phrase of the assistant's compressed answers yet. Together they propose a \emph{dependency map} of the proof (Section~\ref{ch12:sec:theorem}): the steps of the proof and the facts that each step relies on. Once we mark which entries we already know, we call it a \emph{prerequisite map}; its new entries are all we need to learn. This proposed map is incomplete, as Section~\ref{ch17:sec:map} will show. Sections~\ref{ch17:sec:balls}, \ref{ch17:sec:pigeonhole} and~\ref{ch17:sec:subsequence} supply the three missing pieces, each in a short self-contained lesson that we call a \emph{capsule}; the last also sets out the plan of the proof, which Section~\ref{ch17:sec:proof} carries out.

\section{Capsule one: small regions that cover the space}\label{ch17:sec:balls}
The first missing piece is a way to describe a whole metric space, at any desired precision, by finitely many small regions called balls.

\par\noindent\begin{minipage}{\linewidth}
\begin{definition}[Ball]\label{ch17:def:ball}
Let $(X,d)$ be a metric space, let $c\in X$, and let $r>0$. The \emph{ball} with center $c$ and radius $r$ is the set of points of $X$ at distance less than $r$ from $c$:
\[
 B(c,r)=\set{x\in X:d(x,c)<r}.
\]
\end{definition}
\end{minipage}\par

The center belongs to its own ball, because $d(c,c)=0<r$. On the real line, $B(c,r)$ is the open interval $(c-r,c+r)$, because $|x-c|<r$ means $c-r<x<c+r$ (Section~\ref{sec:magnitude}). In the space $[0,1]$, the ball $B(0,1/4)$ is $[0,1/4)$: points outside the space do not count.

The assistant's list did not mention it, but the proof also needs one property of balls. If $x$ and $y$ both lie in $B(c,r)$, then the triangle inequality, with $c$ as the stopping point, and symmetry give
\[
 d(x,y)\le d(x,c)+d(c,y)<r+r=2r.
\]
So any two points of a ball of radius $r$ are less than $2r$ apart. We call this the \emph{diameter bound} for balls. In the proof, it turns ``many terms lie in one small ball'' into ``these terms are close to one another.''

A collection of sets \emph{covers} $X$ if every point of $X$ lies in at least one set of the collection.

\par\noindent\begin{minipage}{\linewidth}
\begin{definition}[Totally bounded]\label{ch17:def:totally-bounded}
A metric space $(X,d)$ is \emph{totally bounded} if, for every $r>0$, there are finitely many points $c_1,\ldots,c_k$ of $X$ such that the balls $B(c_1,r),\ldots,B(c_k,r)$ cover $X$. In symbols:
\[
 \forall r>0\ \exists k\in\NN\ \exists c_1,\ldots,c_k\in X\ \forall x\in X\ \exists i\in\set{1,\ldots,k}:\quad d(x,c_i)<r.
\]
\end{definition}
\end{minipage}\par

Read the quantifiers as a game (Section~\ref{ch02:sec:quantifiers}): a skeptic names a radius $r>0$; you answer with a finite list of centers in the space, which may depend on $r$; the skeptic picks any point $x$ of the space; and you win if $x$ is at distance less than $r$ from one of your centers. A smaller radius may need more centers. The definition does not ask for one finite list that works for every radius at once, and for $[0,1]$ none does.

\subsection*{Two examples on the interval $[0,1]$}
\emph{The interval $[0,1]$ is totally bounded.} Let $r>0$. By the Archimedean property (Section~\ref{ch02:sec:quantifiers}), there is an integer $N>1/r$; it is positive, and $1/N<r$. Use the $N+1$ centers $0,\tfrac1N,\tfrac2N,\ldots,1$, which are spaced $1/N$ apart. A point $x\in[0,1]$ that is a center lies in its own ball; otherwise $j/N<x<(j+1)/N$ for two consecutive centers, and
\[
 d\Bigl(x,\frac jN\Bigr)=x-\frac jN<\frac{j+1}N-\frac jN=\frac1N<r.
\]
So the $N+1$ balls of radius $r$ cover $[0,1]$. For instance, $r=0.3$ allows $N=4$, and the five balls of radius $0.3$ centered at $0$, $0.25$, $0.5$, $0.75$ and $1$ cover $[0,1]$.

\emph{No single finite list works for every radius.} An empty list covers nothing, so fix a list of $k\ge1$ points $c_1,\ldots,c_k$ of $[0,1]$, and take the radius $r=1/(2k)$. The $k+1$ points $0,\tfrac1k,\tfrac2k,\ldots,1$ are at least $1/k=2r$ apart. If the $k$ balls $B(c_i,r)$ covered them, we could assign each of these points to one ball containing it. By the pigeonhole principle (Proposition~\ref{ch07:prop:pigeonhole}), two of the points would share a ball, and by the diameter bound they would be less than $2r$ apart, which is impossible. So the list fails for this radius.

\subsection*{Totally bounded spaces are bounded}
A metric space is \emph{bounded} if there is a real number $M$ such that $d(x,y)\le M$ for all points $x$ and $y$ of the space. Every totally bounded space is bounded. If the space is empty, there is nothing to prove; otherwise, cover it by balls of radius $1$ with centers $c_1,\ldots,c_k$, where $k\ge1$, and let $M$ be the largest of the finitely many distances $d(c_i,c_j)$. If $x$ lies in the ball about $c_i$ and $y$ in the ball about $c_j$, the triangle inequality used twice, together with $d(c_j,y)=d(y,c_j)<1$, gives
\[
 d(x,y)\le d(x,c_i)+d(c_i,c_j)+d(c_j,y)<1+M+1.
\]
So $M+2$ bounds every distance in the space.

The converse is false. On any set $X$, let
\[
 d(x,y)=\begin{cases}0&\text{if }x=y,\\1&\text{if }x\ne y.\end{cases}
\]
This is a metric, the \emph{discrete metric} on $X$. The first two requirements of Definition~\ref{def:metric} hold because $d$ is $0$ exactly when $x=y$, is otherwise $1$, and gives $d(x,y)$ and $d(y,x)$ by the same case. For the triangle inequality $d(x,z)\le d(x,y)+d(y,z)$: if $x=z$, the left side is $0$; if $x\ne z$, the left side is $1$, and the point $y$ cannot equal both $x$ and $z$, so at least one distance on the right is $1$.

On $\NN$ with the discrete metric, every distance is at most $1$, so the space is bounded. But a ball of radius $1/2$ contains only its center, $B(n,1/2)=\set n$, since every other number is at distance $1$ from $n$; so finitely many such balls cannot cover $\NN$: the space is not totally bounded. Boundedness asks for a single bound on all distances; total boundedness asks for a finite description of the space at every scale.

\pauseq{Is the real line, with distance $|x-y|$, totally bounded?}{No. A totally bounded space is bounded, but the real line is not: for any proposed bound $M$, the points $0$ and $|M|+1$ are at distance $|M|+1>M$.}

\section{Capsule two: keeping infinitely many indices}\label{ch17:sec:pigeonhole}
A sequence $x_0,x_1,x_2,\ldots$ has one term for each index $n\in\NN$, so it has infinitely many indices even when its values repeat: $x_n=(-1)^n$ takes only the values $1$ and $-1$, yet the value $1$ occurs at the infinitely many indices $0,2,4,\ldots$. Counting points instead would break the assistant's plan for a constant sequence: no ball contains infinitely many different terms, yet a constant sequence converges at once. So the proof keeps track of sets of \emph{indices}: a ball contains infinitely many terms of a sequence if it contains $x_n$ for infinitely many indices $n$, whether or not these terms are different points.

Recall from Section~\ref{ch03:sec:sets} that a set is \emph{finite} if its members can be written down in a list that ends, and \emph{infinite} otherwise.

\par\noindent\begin{minipage}{\linewidth}
\begin{lemma}[Infinite pigeonhole principle]\label{ch17:lem:pigeonhole}
Let $S$ be an infinite set, and let $A_1,\ldots,A_k$ be finitely many sets whose union contains $S$. Then at least one of the sets $S\cap A_1,\ldots,S\cap A_k$ is infinite.
\end{lemma}
\end{minipage}\par

\begin{proof}
Every member of $S$ lies in at least one of the sets $A_i$, so
\[
 S=(S\cap A_1)\cup(S\cap A_2)\cup\cdots\cup(S\cap A_k).
\]
Suppose, for contradiction, that each of the sets $S\cap A_i$ is finite. Write down the members of $S\cap A_1$, then those of $S\cap A_2$, and so on up to $S\cap A_k$. The result is a list that ends, and every member of $S$ appears in it at least once. Crossing out repetitions, which occur for members lying in several of the sets $A_i$, leaves a list of the members of $S$ that still ends. So $S$ is finite, a contradiction.
\end{proof}

This is the pigeonhole principle of Proposition~\ref{ch07:prop:pigeonhole} for infinitely many objects: if the members of $S$ are placed in finitely many boxes $A_1,\ldots,A_k$, some box holds infinitely many of them. Here, though, the boxes may overlap; an object in two boxes simply appears twice in the list.

For example, take $x_n=(-1)^n$ and the balls $B(1,1/2)$ and $B(-1,1/2)$ on the real line. The indices $n$ with $x_n\in B(1,1/2)$ form the set $A_1$ of even numbers, and those with $x_n\in B(-1,1/2)$ form the set $A_2$ of odd numbers. Their union is $\NN$, and the lemma with $S=\NN$ promises that at least one of them is infinite, without saying which. Here both are.

In the proof of the theorem, $S$ will be the current infinite set of indices and $A_i$ the set of indices $n$ for which $x_n$ lies in the $i$th ball of a finite cover; the lemma then supplies one ball containing $x_n$ for infinitely many $n\in S$.

\section{Capsule three: subsequences, and a plan for the proof}\label{ch17:sec:subsequence}
\par\noindent\begin{minipage}{\linewidth}
\begin{definition}[Subsequence]\label{ch17:def:subsequence}
Let $x_0,x_1,x_2,\ldots$ be a sequence. A \emph{subsequence} of it is a sequence of the form
\[
 x_{n_1},x_{n_2},x_{n_3},\ldots,
 \qquad\text{where }n_1<n_2<n_3<\cdots
\]
are natural numbers.
\end{definition}
\end{minipage}\par

In words: we may skip terms, but the terms we keep must stay in their original order, and no index may be used twice. For $x_n=(-1)^n$, the indices $n_k=2k$ give the subsequence $x_2,x_4,x_6,\ldots$, which is the constant sequence $1,1,1,\ldots$. By contrast, $x_1,x_0,x_3,x_2,\ldots$ is not a subsequence, because its indices do not increase, and neither is $x_0,x_0,x_0,\ldots$, which uses the index $0$ again and again.

Keep the two counters apart: $k$ is the position in the subsequence, while $n_k$ is the position of the same term in the original sequence; in the example, the third term of the subsequence is $x_6$. We number the terms of a subsequence from $k=1$, not $0$. This makes no difference: with $y_k=x_{n_k}$, the definitions of convergence and of Cauchy sequences apply to $y_1,y_2,y_3,\ldots$ as before, with thresholds that are at least $1$.

To make the chosen indices increase, we need one fact. \emph{An infinite set $S$ of natural numbers contains a number larger than any given natural number $M$.} Otherwise every member of $S$ would lie in $\set{0,1,\ldots,M}$, and crossing out of the list $0,1,\ldots,M$ every number not in $S$ would leave a list of the members of $S$ that ends, so $S$ would be finite. So if each stage of a construction offers an infinite set of allowed indices, we can always choose an allowed index larger than the one chosen at the previous stage.

\subsection*{A plan with three kinds of objects}
The construction involves three kinds of objects that must not be mixed up: points of the space, such as the term $x_7$; balls, which are sets of points; and sets of indices, which are sets of natural numbers. If $x_7=x_9$ lies in a chosen ball, it is one point of the ball but contributes two indices, $7$ and $9$.

Here is the plan. Start with every index allowed: $S_0=\NN$. At stage~$1$, cover the space by finitely many balls of radius $1/4$, as total boundedness allows. By Lemma~\ref{ch17:lem:pigeonhole}, one of these balls contains $x_n$ for infinitely many indices $n\in S_0$; keep exactly those indices, as the set $S_1$. By the diameter bound, any two terms with indices in $S_1$ are less than $2\cdot\tfrac14=\tfrac12$ apart. In general, stage~$k$ uses balls of radius $2^{-(k+1)}$, considers only the indices kept at stage~$k-1$, and keeps infinitely many indices whose terms are less than $2\cdot2^{-(k+1)}=2^{-k}$ apart. The radius is half the desired distance because the diameter bound doubles it; any radii shrinking towards zero would work, and this choice keeps the arithmetic tidy. Since each stage only throws indices away,
\[
 S_0\supseteq S_1\supseteq S_2\supseteq\cdots,
\]
where $A\supseteq B$, read ``$A$ contains $B$,'' means $B\subseteq A$. We say that the index sets are \emph{nested}. The chosen balls need not be nested: the ball of stage~$2$ may stick out of the ball of stage~$1$. That does no harm, because a term whose index lies in $S_2$ also has its index in $S_1$, so it lies in both chosen balls. The nesting of the index sets carries the distance bound of each stage forward to all later stages.

\section{The complete proof}\label{ch17:sec:proof}
\par\noindent\begin{minipage}{\linewidth}
\begin{theorem}\label{thm:subsequence}
Let $(X,d)$ be a complete, totally bounded metric space. Then every sequence $x_0,x_1,x_2,\ldots$ in $X$ has a subsequence that converges to a point of $X$.
\end{theorem}
\end{minipage}\par

Unlike the contraction theorem (Theorem~\ref{thm:banach}), this theorem does not assume that $X$ is nonempty: if $X$ has no points, there is no sequence in $X$ at all, and the statement is vacuously true (Section~\ref{ch02:sec:contrapositive}).

\begin{proof}
Let $x_0,x_1,x_2,\ldots$ be a sequence in $X$. The proof has five steps.

\emph{Step 1: nested infinite sets of indices.} We construct, one stage at a time, infinite sets of indices $S_0\supseteq S_1\supseteq S_2\supseteq\cdots$ such that, for every $k\ge1$,
\begin{equation}\label{ch17:eq:stage}
 d(x_i,x_j)<2^{-k}\qquad\text{for all }i,j\in S_k.
\end{equation}
Start with $S_0=\NN$, which is infinite. Now let $k\ge1$, and suppose that $S_{k-1}$ has already been constructed and is infinite. Since $X$ is totally bounded, finitely many balls of radius $2^{-(k+1)}$ cover $X$. For each of these balls, consider the set of indices $n$ for which $x_n$ lies in that ball. Every term $x_n$ lies in at least one of the balls, so every index lies in at least one of these finitely many sets. By Lemma~\ref{ch17:lem:pigeonhole}, applied to the infinite set $S_{k-1}$, there is a ball $B$ among them such that the set
\[
 S_k=\set{n\in S_{k-1}:x_n\in B}
\]
is infinite. By its definition, $S_k\subseteq S_{k-1}$. If $i,j\in S_k$, then $x_i$ and $x_j$ both lie in the ball $B$ of radius $2^{-(k+1)}$, so the diameter bound gives $d(x_i,x_j)<2\cdot2^{-(k+1)}=2^{-k}$. This proves~\eqref{ch17:eq:stage} for this value of $k$. Since $S_k$ is again infinite, the next stage can begin, and the construction continues for $k=1,2,3,\ldots$.

\emph{Step 2: increasing indices.} Choose any $n_1\in S_1$. For $k\ge2$, once $n_{k-1}$ has been chosen, choose $n_k\in S_k$ with $n_k>n_{k-1}$. Such an index exists because $S_k$ is infinite, and an infinite set of natural numbers contains a number larger than any given one (Section~\ref{ch17:sec:subsequence}). Then $n_1<n_2<n_3<\cdots$, so the terms $x_{n_1},x_{n_2},x_{n_3},\ldots$ form a subsequence. By construction, $n_k\in S_k$ for every $k\ge1$.

\emph{Step 3: later terms keep the earlier bounds.} Let $1\le k\le\ell$. The sets are nested, so $S_\ell\subseteq S_{\ell-1}\subseteq\cdots\subseteq S_k$, and therefore $n_\ell\in S_\ell\subseteq S_k$. So not only the term chosen at stage~$k$, but every term chosen at a later stage, has its index in $S_k$. In particular, if $\ell\ge k$ and $m\ge k$, then $n_\ell$ and $n_m$ both lie in $S_k$, and~\eqref{ch17:eq:stage} gives
\[
 d(x_{n_\ell},x_{n_m})<2^{-k}.
\]

\emph{Step 4: the subsequence is a Cauchy sequence.} Let $\eps>0$. By the Archimedean property (Section~\ref{ch02:sec:quantifiers}), choose an integer $K$ with $K>1/\eps$. It is positive, because $1/\eps>0$, so $K\ge1$. This matters, because the bound~\eqref{ch17:eq:stage} holds only from stage~$1$ on; the set $S_0=\NN$ carries no information about distances. Section~\ref{ch11:sec:convergence} showed that $2^K\ge K+1$, so
\[
 2^{-K}\le\frac1{K+1}<\frac1K<\eps.
\]
For any $\ell,m\ge K$, Step~3 with $k=K$ gives $d(x_{n_\ell},x_{n_m})<2^{-K}<\eps$. So $K$ is a threshold that works for the tolerance $\eps$. Since $\eps>0$ was arbitrary, the subsequence is a Cauchy sequence.

\emph{Step 5: completeness supplies the limit.} The space $X$ is complete, so the Cauchy sequence $x_{n_1},x_{n_2},x_{n_3},\ldots$ converges to some point $x_*$ of $X$. This is the required subsequence.
\end{proof}

\subsection*{Watching the construction on an example}
The interval $[-1,2]$, with distance $|x-y|$, is complete, because it is a closed interval (Section~\ref{ch11:sec:complete}). It is totally bounded by the same argument as for $[0,1]$ in Section~\ref{ch17:sec:balls}; Exercise~\ref{ex:17.4} asks you to write out the details for every interval $[a,b]$ with $a<b$. In this space consider the sequence
\[
 x_n=(-1)^n+\frac1{n+1},
\]
whose first terms are $x_0=2$, $x_1=-\tfrac12$, $x_2=\tfrac43$, $x_3=-\tfrac34$ and $x_4=\tfrac65$. The even-indexed terms $1+\tfrac1{n+1}$ crowd towards $1$ from above, and the odd-indexed terms $-1+\tfrac1{n+1}$ crowd towards $-1$, never rising above $-\tfrac12$.

At stage~$1$, cover $[-1,2]$ by the thirteen balls of radius $1/4$ centered at $-1,-\tfrac34,-\tfrac12,\ldots,\tfrac74,2$. These balls do cover the interval: a point of $[-1,2]$ that is not a center lies strictly between two consecutive centers, which are $1/4$ apart, so its distance from the left one is less than $1/4$. The ball $B(1,1/4)$ is one of the thirteen. It contains $x_n$ exactly when $n$ is even and $\tfrac1{n+1}<\tfrac14$, that is, when $n$ is even and $n\ge4$. So the construction may take $S_1=\set{4,6,8,\ldots}$. Other balls also contain infinitely many terms, for instance $B(-1,1/4)$, which contains $x_n$ for every odd $n\ge5$, and the construction is free to pick any of them.

At stage~$2$, use the balls of radius $1/8$ centered at the points $-1,-\tfrac78,\ldots,2$, which include $1$. Among the indices in $S_1$, the ball $B(1,1/8)$ contains $x_n$ exactly when $\tfrac1{n+1}<\tfrac18$, that is, when $n\ge8$. So we may take $S_2=\set{8,10,12,\ldots}$, which lies inside $S_1$. In the same way, stage~$k$ can use the balls of radius $2^{-(k+1)}$ centered at the points $-1+j\cdot2^{-(k+1)}$ for $j=0,1,\ldots,3\cdot2^{k+1}$, and $1$ is always one of these centers: $j=2^{k+2}$ gives $-1+2^{k+2}\cdot2^{-(k+1)}=-1+2=1$. If we keep choosing the ball around $1$, stage~$k$ gives the set $S_k$ of even numbers that are at least $2^{k+1}$, because $\tfrac1{n+1}<2^{-(k+1)}$ means $n+1>2^{k+1}$, that is, $n\ge2^{k+1}$.

Now choose the indices of Step~2, taking each time the smallest allowed index. This gives $n_1=4$, then $n_2=8$ (the smallest member of $S_2$ larger than $4$), then $n_3=16$, and in general $n_k=2^{k+1}$. The subsequence is $x_4,x_8,x_{16},\ldots$, with terms $1+\tfrac1{2^{k+1}+1}$, and it converges to $1$. Had stage~$1$ chosen the ball $B(-1,1/4)$ instead, every later index would have been odd, and the construction would have produced a subsequence converging to $-1$. The theorem promises \emph{some} convergent subsequence, not a particular one. In this example the indices follow the simple formula $n_k=2^{k+1}$; in general there is no such simple formula. What matters is only that each $n_k$ lies in $S_k$ and is larger than $n_{k-1}$.

\subsection*{Why the proof picks a new index at every stage}
It is tempting to look for a single index $n$ that lies in every set $S_k$, so that the one term $x_n$ lies in every chosen ball. In general there is no such index, because nested infinite sets can have no member in common. For example, the sets $\set{k,k+1,k+2,\ldots}$ for $k=1,2,3,\ldots$ are infinite and nested, but no natural number $n$ belongs to all of them, since $n$ is not a member of $\set{n+1,n+2,\ldots}$. The same happens in the construction we just watched. There $S_k$ contains only numbers that are at least $2^{k+1}$, and a natural number $n$ is smaller than $2^{k+1}$ as soon as $k\ge n$, because $2^{k+1}>2^k\ge k+1>n$. So no index lies in every $S_k$. Even when such an index exists, as it does for a constant sequence, it gives only one term, not a subsequence. This is why the proof picks a new, later index from each set in turn. Step~3 shows that any two sufficiently late picks lie together in one earlier set, and that is exactly what the Cauchy condition asks for.

\pauseq{Where does the proof use each hypothesis of Theorem~\ref{thm:subsequence}?}{There are two hypotheses. Total boundedness is used at every stage of Step~1, to cover $X$ by finitely many balls of radius $2^{-(k+1)}$. The finiteness of that cover is what allows Lemma~\ref{ch17:lem:pigeonhole} to keep infinitely many indices. Completeness is used once, in Step~5, to give the Cauchy subsequence a limit in $X$. The other ingredients, namely the diameter bound, the nesting of the index sets and the increasing choice of indices, are not hypotheses. They are facts proved along the way.}

\section{Why neither hypothesis can be dropped}\label{ch17:sec:counterexamples}
As in Section~\ref{ch12:sec:hypotheses}, we remove one hypothesis at a time, keep the other, and look for an example in which the conclusion fails.

\subsection*{Without total boundedness}
Give $\NN$ the discrete metric of Section~\ref{ch17:sec:balls}, in which different numbers are at distance $1$. This space is complete. To see why, take a Cauchy sequence in it and apply the Cauchy condition with the tolerance $1/2$. It supplies a threshold $N$ such that any two terms from $x_N$ onward are at distance less than $1/2$. Distances in this space are only ever $0$ or $1$, so all these distances are $0$, and every term from $x_N$ onward equals $x_N$. Such a sequence converges to $x_N$, a point of the space: for every $\eps>0$ and every $n\ge N$, we have $d(x_n,x_N)=0<\eps$.

The space is not totally bounded, as Section~\ref{ch17:sec:balls} showed. And the conclusion of the theorem fails. Consider the sequence $x_n=n$, that is, $0,1,2,\ldots$. In any subsequence $x_{n_1},x_{n_2},\ldots$ the indices increase, so the terms $n_1,n_2,\ldots$ are different numbers, and any two of them are at distance $1$. For the tolerance $1/2$, no threshold can work, so the subsequence is not a Cauchy sequence. Since every convergent sequence is a Cauchy sequence (Proposition~\ref{ch11:prop:convergent-cauchy}), the subsequence does not converge. In the proof, it is Step~1 that breaks down: no finite collection of balls of radius $1/4$ covers the space, so the selection cannot even begin.

\subsection*{Without completeness}
The interval $(0,1)$, with distance $|x-y|$, is totally bounded. The centers $0$ and $1$ used for $[0,1]$ are not allowed here, because the centers must be points of the space, so we use midpoints instead. Given $r>0$, the Archimedean property supplies an integer $N>1/(2r)$, which is positive and satisfies $1/(2N)<r$. Split $[0,1]$ into the $N$ pieces $\bigl[\tfrac{j-1}N,\tfrac jN\bigr]$ for $j=1,\ldots,N$. Their midpoints $\tfrac{2j-1}{2N}$ lie strictly between $0$ and $1$, so they are points of $(0,1)$. Every $x\in(0,1)$ lies in one of the pieces. A point of a piece is at most half the length of the piece away from its midpoint, so $x$ is at distance at most $1/(2N)<r$ from that midpoint.

But $(0,1)$ is not complete (Section~\ref{ch11:sec:complete}), and the conclusion of the theorem fails. Consider the sequence $x_n=1/(n+2)$, that is, $\tfrac12,\tfrac13,\tfrac14,\ldots$, whose terms all lie in $(0,1)$. On the real line it converges to $0$: given $\eps>0$, choose an integer $N>1/\eps$; then every $n\ge N$ gives $x_n=\tfrac1{n+2}<\tfrac1N<\eps$.

Every subsequence $x_{n_1},x_{n_2},\ldots$ also converges to $0$ on the real line. The indices increase by at least one at every step, starting from $n_1\ge0$, so induction gives $n_k\ge k-1$ for every $k\ge1$. Hence, with $\eps$ and $N$ as above, every $k\ge N+1$ has $n_k\ge N$ and therefore $x_{n_k}<\eps$. Now suppose some subsequence converged to a point $L$ of $(0,1)$. Distances in $(0,1)$ are distances on the real line, so the subsequence would converge to $L$ on the real line too. It also converges to $0$ there, so uniqueness of limits (Proposition~\ref{thm:limitunique}) gives $L=0$, which is not a point of $(0,1)$. This contradiction shows that no subsequence converges in $(0,1)$. In the proof, Steps~1--4 go through, and it is Step~5 that fails: the Cauchy subsequence exists, but the point it crowds around is missing from the space.

Together, the two examples show that the hypotheses do different jobs: total boundedness produces a Cauchy subsequence, and completeness gives it a limit. When you meet a new hypothesis, ask not only ``What does it mean?'' but also ``Which step fails without it?''

\section{The prerequisite map after the work is done}\label{ch17:sec:map}
The assistant's list of prerequisites in Section~\ref{ch17:sec:start} was incomplete. It did not mention the diameter bound for balls, used in Step~1, or explain why nested selections make the chosen terms a Cauchy sequence, the argument of Step~3. Such gaps turn up when you rebuild a proof line by line, and finding them is part of the method: add each missing step to the map, learn it, and return to the argument. A list of prerequisites deserves your trust only once you have used it to rebuild the proof.

Here is the completed map.
\begin{center}
\begin{tabularx}{\linewidth}{@{}l>{\raggedright\arraybackslash}X>{\raggedright\arraybackslash}X@{}}
\toprule
Where & Facts used & Conclusion\\
\midrule
Capsule one & triangle inequality and symmetry & diameter bound for balls\\
Capsule two & definition of a finite set & infinite pigeonhole principle (Lemma~\ref{ch17:lem:pigeonhole})\\
Capsule three & definition of a finite set & an infinite set of natural numbers contains a number larger than any given one\\
Step~1 & total boundedness, infinite pigeonhole principle, diameter bound & nested infinite index sets $S_k$ satisfying~\eqref{ch17:eq:stage}\\
Step~2 & the sets $S_k$ are infinite, and an infinite set of natural numbers contains a number larger than any given one & increasing indices $n_k\in S_k$, hence a subsequence\\
Step~3 & Step~2 (each $n_\ell$ lies in $S_\ell$), nesting of the sets $S_k$, the bound~\eqref{ch17:eq:stage} & terms from position $k$ on are less than $2^{-k}$ apart\\
Step~4 & Step~3, Archimedean property, $2^K\ge K+1$ & a Cauchy subsequence\\
Step~5 & Step~4, completeness & a subsequence that converges to a point of the space\\
\bottomrule
\end{tabularx}
\end{center}
The facts not proved here came from Chapter~\ref{ch:limits} (the triangle inequality, Cauchy sequences, completeness, $2^K\ge K+1$), Chapter~\ref{ch:logic} (the Archimedean property) and Chapter~\ref{ch:maps} (finite sets). Each new ingredient was put to work in the proof soon after it was learned. That matters: a concept you have seen doing its job stays with you, while one learned and left unused becomes disconnected vocabulary.

\subsection*{Check what you have learned, not only what you can name}
Here is a four-part check of whether you have learned a new concept well enough for the proof at hand, carried out for total boundedness.
\begin{enumerate}
\item \emph{State the problem the concept solves.} At every scale, the proof must split the infinitely many indices into finitely many groups whose terms are close together. Total boundedness provides, for any radius, finitely many balls of that radius covering the space, and the balls give the groups.
\item \emph{State the definition with its quantifiers in the right order.} The radius comes first. The list of centers may depend on the radius, but for that radius one list must cover every point.
\item \emph{Test it on examples near the boundary.} The discrete metric on $\NN$ shows that a bound on all distances does not give finite covers by small balls. The interval $(0,1)$ shows that finite covers by small balls do not provide the missing limits.
\item \emph{Return to the exact line of the proof where the concept is used,} here the beginning of each stage of Step~1, and explain it. Why must the cover be finite? Because Lemma~\ref{ch17:lem:pigeonhole} applies only to finitely many sets: with infinitely many balls, each might contain just one term, as in the discrete metric on $\NN$, and no ball would keep infinitely many indices.
\end{enumerate}
If you can recite the definition but cannot explain that line, you have not yet learned the concept well enough for this proof.

\subsection*{Labels for the entries of the map}
For an unfamiliar paper, mark each entry of the map with its status, using the labels of ``Keeping track of evidence'' in \emph{To the reader} and one new label, \emph{borrowed}.
\begin{itemize}
\item \emph{Understood:} the entry passes the four-part check above, and you can use it in a slightly different setting. After studying $[0,1]$, proving that every interval $[a,b]$ is totally bounded (Exercise~\ref{ex:17.4}) is a small test of this kind.
\item \emph{Can follow:} you can check each step but could not yet explain the entry, or use it in a new setting, on your own. You might check every line of the proof of Lemma~\ref{ch17:lem:pigeonhole} and still be unable to say, without looking, which line would fail with infinitely many sets $A_i$.
\item \emph{Borrowed:} you use the result without having studied its proof. You must still check that its hypotheses hold and that its conclusion is exactly what you need; learning its proof can be a separate task. Here the completeness of the real line is borrowed: Chapter~\ref{ch:limits} accepted it as a starting fact, and the completeness of $[-1,2]$ in Section~\ref{ch17:sec:proof} rests on it.
\item \emph{Open:} you cannot yet follow the entry, which records exactly what remains to be learned, like much of the research record of Section~\ref{ch17:sec:research}.
\end{itemize}
A map may contain borrowed results and open entries, as long as it labels them clearly.

\section{Reading a formalized research proof with the same method}\label{ch17:sec:research}
We now apply the same method to the two research claims met earlier in the book: Seymour's second-neighborhood conjecture (Section~\ref{ch07:sec:seymour}) and the circulant Hadamard conjecture (Section~\ref{ch14:sec:research}). Claimed proofs of both appear in OpenAI's public mathematics repository~\cite{oairepo}. The repository collects preprints, in the sense of Section~\ref{ch07:sec:seymour}, together with formal proofs in Lean. For each problem it also describes exactly which statements were formalized. That a paper appears in the repository does not, by itself, vouch for every claim in the paper. The right unit of reading is one exact theorem, together with its definitions and the results it depends on, not a title or an impressive final line.

\subsection*{Decoding the statement}
Start with Seymour's conjecture. The repository's formal statement is about a nonempty finite oriented graph, and it asks for a vertex whose second out-neighborhood has at least as many vertices as its first~\cite{oai173}. Section~\ref{ch07:sec:seymour} supplies the mathematical vocabulary, and Section~\ref{ch16:sec:second} showed how the exclusions in the definition of a second out-neighbor are written in Lean. The official statement lives in a file of its own~\cite{oaistatement}, and the theorem is called \code{exists\_goodVertex}. That file declares the assumptions on separate lines and uses the symbols $\to$ and $\exists$. Collecting the assumptions in one place, writing \code{->} and \code{Exists} as in Chapter~\ref{ch:lean2}, and putting \code{Type} where the file has \code{Type*} (a technical detail about the sizes of types that makes no difference here), we can display the statement as follows.
\par\noindent\begin{minipage}{\linewidth}
\begin{lstlisting}
theorem exists_goodVertex {V : Type}
    [Fintype V] [DecidableEq V] [Nonempty V]
    (r : V -> V -> Prop) (hr : IsOriented r) :
    Exists (fun v => GoodVertex r v) := by
  sorry
\end{lstlisting}
\end{minipage}\par
The repository builds on Mathlib, the large mathematical library mentioned in Chapter~\ref{ch:lean1}, which provides finite sets and ways of counting their members. Reading the statement piece by piece, we find the following.
\begin{itemize}
\item \code{V} is the type of vertices. Three assumptions are written in square brackets. \code{[Fintype V]} says that there are only finitely many vertices, and \code{[Nonempty V]} says that there is at least one. \code{[DecidableEq V]} says that, for any two vertices, it can be decided whether they are equal. This last assumption is a technical requirement of Mathlib's finite sets, and it does not exclude any graph: with the standard axiom \code{Classical.choice} of Section~\ref{ch16:sec:axioms}, Lean can supply it for every type.
\item The edge relation \code{r} has type \code{V -> V -> Prop}, so \code{r v w} is the proposition $v\to w$, exactly as in Section~\ref{ch16:sec:second}.
\item The explicit hypothesis \code{(hr : IsOriented r)} says that \code{r} has no loops and no pair of opposite edges. The definition of \code{IsOriented} has exactly these two parts, named \code{loopless} and \code{asymmetric}.
\item The conclusion says that some vertex \code{v} satisfies \code{GoodVertex r v}. This is defined to mean that the set \code{firstNeighbors r v} has at most as many members as the set \code{secondNeighbors r v}. The first of these sets consists of the vertices \code{w} with \code{r v w}, that is, the out-neighbors of $v$. The second is defined by exactly the three conditions of Section~\ref{ch07:sec:seymour}: $w\ne v$, $v\to w$ fails, and $v\to u\to w$ for some vertex $u$.
\end{itemize}
Each piece of the formal statement matches a piece of the informal one: ``nonempty'' and ``finite'' are the bracketed assumptions, ``oriented'' is \code{IsOriented}, and the inequality $|N_2^+(v)|\ge|N_1^+(v)|$ is \code{GoodVertex}. No condition of the informal statement is missing. Reading the statement in this way needed nothing beyond Chapters~\ref{ch:graphs}, \ref{ch:lean1} and~\ref{ch:lean2}. After Section~\ref{ch16:sec:axioms}, the word \code{sorry} at the end may look alarming. It belongs there: this file only records what must be proved, and the proof is in another file.

\subsection*{Locating the main steps}
The proof is in a different file~\cite{oaimain}. That file imports three further files, named \code{Reduction}, \code{Pruning} and \code{ExtremalContradiction}, and its own proof is only a few lines long. It argues by contradiction. It assumes that no vertex is good, so that the graph is a counterexample; it applies a theorem from the reduction file to turn the counterexample into a new oriented graph with additional properties; and it then applies results from the other two files to reach a contradiction.

The reduction theorem~\cite{oaireduction} says that from any counterexample one can build, for some positive integer $n$, an oriented graph whose vertices are the pairs $(i,j)$ of natural numbers with $i<n$ and $j\le n$, and which has two further properties, named \code{PositiveIndegree} and \code{StrictSubsetGrowth}. The reduction file in turn imports three more files, so the map keeps branching. The name \code{PositiveIndegree} suggests that every vertex has at least one edge coming in, but a name is not a definition: as Section~\ref{ch15:sec:symbols} showed, Lean accepts any name for any statement. To understand these properties, you have to open their definitions and read them.

\subsection*{Keeping a reading record}
A \emph{reading record} is a written account of a proof you are studying: what you have read, what you have understood, and what remains open. Section~\ref{appE:sec:record} describes what a finished one contains. At this point, a reading record of the Seymour proof has entries of three different kinds.
\begin{itemize}
\item \emph{Understood:} the final statement has been decoded, and its definitions match the informal conjecture.
\item \emph{Borrowed:} the logical role of the reduction is clear, since it turns a counterexample into a graph with extra properties, but its proof has not been read.
\item \emph{Open:} the meaning of \code{PositiveIndegree} and \code{StrictSubsetGrowth}, and the mathematics of the pruning step and of the final contradiction, remain to be studied from their sources.
\end{itemize}
The proof as a whole is long and uses ideas well beyond this course. A dependency sketch like this one is a way into the proof, not a replacement for it.

\subsection*{What the formal check compares}
How does the repository connect the official statement to the proof? A configuration file~\cite{oaiconfig} lists what the formal check compares: the file containing the official statement, the file containing the proof, the name of the theorem, and the axioms the proof may use. The permitted axioms are \code{propext}, \code{Quot.sound} and \code{Classical.choice}, the three standard axioms of Section~\ref{ch16:sec:axioms}. The check is designed to confirm that the proof file proves exactly the official statement, using no axioms beyond the permitted ones. It is the two statements that are compared, not merely their names: a theorem in the proof file with the same name as the official one would, by itself, say nothing about whether the two statements agree.

Reading the configuration file tells you what the check is designed to compare; it does not tell you that the check passed. Unless you run the check yourself, the statement that it passed is a claim made by the repository, and your reading record should attribute it: ``the repository reports that the check passes,'' rather than ``the check passes.'' In the same way, the record can describe the role of a file whose proofs you have not read, but it cannot vouch for those proofs. What it can report on your own authority is what you decoded or reconstructed yourself.

\subsection*{The circulant Hadamard conjecture}
For the circulant Hadamard conjecture of Chapter~\ref{ch:signs}, Section~\ref{ch14:sec:research} has already read the repository's description of exactly which statements were formalized~\cite{oai179}. You have also worked through the proof of a necessary condition yourself: by Proposition~\ref{thm:circulantnecessary}, a circulant Hadamard matrix of order greater than $4$ would need an order $4u^2$ with $u\ge2$, such as $16$, $36$ or $64$. That gives a precise next question: which lemma of the proof rules out these orders? A question informed by mathematics you have done yourself will get you much further than asking an assistant to ``explain the whole repository.''

\section*{Exercises}
Exercises~\ref{ex:17.1}--\ref{ex:17.4} are core exercises, and Exercises~\ref{ex:17.5} and~\ref{ex:17.6} are transfer exercises. Complete solutions are in Appendix~\ref{app:solc}, under ``Week~17.''
\begin{exercise}\label{ex:17.1}Suppose a ball contains $x_n$ for infinitely many indices $n$. Explain why it does not follow that the sequence takes infinitely many different values in that ball. Does the proof of Theorem~\ref{thm:subsequence} need the sequence to take infinitely many different values?\end{exercise}
\begin{exercise}\label{ex:17.2}An assistant writing out the proof of Theorem~\ref{thm:subsequence} chooses an arbitrary $n_k\in S_k$ at each stage, without requiring $n_k>n_{k-1}$. Identify the gap and repair the choice. Explain why the repair is always possible.\end{exercise}
\begin{exercise}\label{ex:17.3}Let $X$ be any set with the discrete metric of Section~\ref{ch17:sec:balls}. (a)~Prove that $X$ is complete. (b)~Prove that $X$ is totally bounded if and only if $X$ is finite. (c)~Prove that a sequence $x_0,x_1,x_2,\ldots$ in $X$ has a subsequence that converges in $X$ exactly when there is a point $c$ of $X$ with $x_n=c$ for infinitely many indices $n$. Use this to explain, for $X=\NN$ and the sequence $x_n=n$, the difference between completeness and the conclusion of Theorem~\ref{thm:subsequence}.\end{exercise}
\begin{exercise}\label{ex:17.4}Let $a<b$ be real numbers. Give a complete proof that the interval $[a,b]$, with distance $|x-y|$, is totally bounded. State exactly what the number of balls in your cover may depend on, and what it may not depend on.\end{exercise}
\begin{exercise}\label{ex:17.5}A directed graph has the three vertices $v$, $u$, $w$ and the edges $v\to u$, $u\to w$ and $v\to w$. Let $R_2(v)$ be the set of vertices $z$ for which there is a route $v\to y\to z$ through some vertex $y$. Find $N_1^+(v)$, $N_2^+(v)$ and $R_2(v)$. Then explain what would go wrong if a formal statement of Seymour's conjecture used $R_2(v)$ in place of $N_2^+(v)$.\end{exercise}
\begin{exercise}\label{ex:17.6}Write a minimal-background request, in the sense of Section~\ref{ch17:sec:start}, for the proof of the contraction theorem (Theorem~\ref{thm:banach}), and another for the proof of Theorem~\ref{thm:subsequence}. For each of the two proofs, name a concept that it needs but the other proof does not. Then state what a reading record of a formalized research proof may honestly claim if you have not run its formal check yourself.\end{exercise}

